\documentclass[11pt]{article}

\usepackage[margin=1in]{geometry}
\usepackage{amsmath, amssymb, amsthm}
\usepackage{bm}
\usepackage{hyperref}
\usepackage{xcolor}

\hypersetup{colorlinks=true, linkcolor=blue!50!black, urlcolor=blue!50!black}

% --- notation ---------------------------------------------------------------
\newcommand{\E}{\mathbb{E}}
\newcommand{\Eq}[1][q]{\mathbb{E}_{#1}}
\newcommand{\Eqx}{\mathbb{E}_{q(x_i)}}
\newcommand{\Eqxx}{\mathbb{E}_{q(x_i,\,x_{i+1})}}
\newcommand{\Cov}{\operatorname{Cov}}
\newcommand{\diag}{\operatorname{diag}}
\newcommand{\tr}{\operatorname{tr}}
\newcommand{\GP}{\operatorname{GP}}
\newcommand{\given}{\,|\,}

\newcommand{\code}[1]{\texttt{\small #1}}
% Long unbreakable \code{...} identifiers otherwise push prose into the margin.
\emergencystretch=3em
\hbadness=10000

% Section banner for ELBO-term grounding.  Use immediately under \section{...}.
\newcommand{\elbobanner}[2]{%
\par\noindent\textbf{\textit{ELBO term:}}\ #1\quad\textbf{\textit{Quantity produced:}}\ #2\par\medskip}

\title{The EFGP-SING E-step\\
\large GP regression on uncertain latent velocities,
solved by Fourier features}
\author{}
\date{}

\begin{document}
\maketitle

\begin{abstract}
\noindent
This note specifies the E-step of SING with a GP drift prior,
accelerated by EFGP. \S\ref{sec:interface} states the whole contract:
SING's natural-gradient backbone asks a drift model for exactly three
per-transition moments and \code{jax.grad}s through them, so EFGP is
fully described by three values plus six VJP slots, together with the
three departures (D1--D3) from the generic SING algorithm. The rest is
the \emph{why}: the $q(f)$ update is GP regression on latent velocities
(\S\ref{sec:qf-update}), made matrix-free by Stein's lemma plus a
Fourier basis whose Gram is BTTB Toeplitz (\S\ref{sec:efgp}); the
remaining $q(x)$-expectations are Gaussian-enveloped NUFFTs, computed by
one of three interchangeable backends (\S\ref{sec:compute}); and the
$q(f)\!\to\!q(x)$ handoff is closed by a Gauss--Newton truncation of
Price's theorem, which is what makes the load-bearing $\partial/\partial
S_i$ slot PSD by construction (\S\ref{sec:qx}).
\end{abstract}

\tableofcontents

% ============================================================================
\section{Setup and ELBO road-map}
\label{sec:setup}
% ============================================================================

\subsection*{Latent SDE and observations}

We assume an Euler-discretised latent SDE
\begin{equation}
x_{i+1} \;=\; x_i \;+\; \Delta_i\, f(x_i) \;+\; \varepsilon_i,
\qquad \varepsilon_i \sim \mathcal{N}(0,\,\Delta_i\Sigma),
\label{eq:sde}
\end{equation}
with Gaussian observations $p(y_i\given x_i)$. Throughout, we abbreviate
\begin{equation}
\delta_i \;:=\; x_{i+1} - x_i.
\end{equation}
The diffusion is taken diagonal,
$\Sigma=\diag(\sigma_1^2,\ldots,\sigma_{D_x}^2)$. The drift coordinates carry
independent GP priors,
\begin{equation}
f_r \sim \GP(0, k_\theta), \qquad r = 1,\ldots,D_x.
\label{eq:gp-prior}
\end{equation}

\subsection*{Joint, ELBO, and the three vEM updates}

The joint factorises as
\begin{equation}
p(f, x_{0:T}, y_{0:T})
\;=\;
p(f)\,p(x_0)
\prod_i p(x_{i+1}\given x_i, f)
\prod_i p(y_i\given x_i),
\end{equation}
with structured variational posterior $q(x_{0:T})\,q(f)$. The ELBO is
\begin{equation}
\begin{aligned}
\mathcal{L}(q)
\;=\;
&\underbrace{\E_{q(x)}\!\bigl[\log p(y\given x)\bigr]}_{\substack{\text{Gaussian closed-form;}\\\text{$q(x)$-grad: backbone autograd}}}
\;+\;
\underbrace{\sum_i \E_{q(x,f)}\!\bigl[\log p(x_{i+1}\given x_i, f)\bigr]}_{\substack{\text{transition --- couples $q(x), q(f)$;}\\\text{$q(x)$-grad: \emph{hand-derived} \code{custom\_vjp} (§\ref{sec:qx});}\\\text{$q(f)$-update: closed-form $A_r\mu_r{=}h_r$ (§\ref{sec:qf-update})}}}
\\[1.2em]
&\;+\;
\underbrace{\E_{q(f)}\!\bigl[\log p(f) - \log q(f)\bigr]}_{\substack{\text{constant in $q(x)$;}\\\text{enters $q(f)$-update + M-step}}}
\;+\;
\underbrace{H[q(x)] + \E_{q(x_0)}\!\bigl[\log p(x_0)\bigr]}_{\substack{\text{Gaussian closed-form;}\\\text{$q(x)$-grad: backbone autograd}}}.
\end{aligned}
\label{eq:elbo}
\end{equation}

vEM cycles three coupled updates against this ELBO:

\begin{itemize}
\setlength\itemsep{2pt}
\item[\textbf{q(f) update}\ ] (§\ref{sec:qf-update}–§\ref{sec:compute}):
maximise $\mathcal{L}$ over $q(f)$ with $q(x)$ fixed. Closed-form Gaussian
$q(w_r)=\mathcal{N}(\mu_r, A_r^{-1})$; $A_r\mu_r=h_r$ via Jacobi-PCG.

\item[\textbf{q(x) update}\ ] (§\ref{sec:qx}):
SING $\rho$-damped natural gradient on $\mathcal{L}$ with $q(f)$ fixed.
The transition contribution's $q(x)$-integrals are closed in $x$ by
Stein (for the pairwise piece) and a first-order Taylor linearisation
at $m_i$ (for the residual nonlinearity in $\bar f$ and $V$).

\item[\textbf{M-step}\ ] (separate doc): kernel hypers $\theta$ maximise the
collapsed marginal NLL. At the q(f)-optimum,
\begin{equation}
\mathcal{L}_{\rm M}(\theta)
\;=\;
-\frac{1}{2}\sum_r \mathrm{Re}\bigl\langle h_{\theta,r},\, \mu_{\theta,r}\bigr\rangle
\;+\;\frac{D_{\rm out}}{2}\,\log\bigl|A_\theta\bigr|.
\label{eq:Mstep-collapsed}
\end{equation}
\end{itemize}

Only the transition term in \eqref{eq:elbo} couples $q(x)$ and $q(f)$
non-trivially; the rest of this document analyses that term.

\subsection*{Section map}

\begin{center}
\begin{tabular}{@{}l l p{0.50\textwidth}@{}}
\hline
\S & Addresses & Produces \\
\hline
\ref{sec:interface} & \textbf{the SING contract} &
  \textbf{3 required moments, 6 VJP slots, departures D1--D3} \\
\ref{sec:notation} & — & Notation glossary \\
\ref{sec:qf-update} (incl.\ \ref{sec:transition}–\ref{sec:uncertain-x}) & transition at fixed $q(x)$ &
  $A_r, h_r$ (basis-agnostic) \\
\ref{sec:efgp} & same & Fourier specialisation: BTTB + cheap Jacobian \\
\ref{sec:compute} & same &
  the $q(x)$-expectations that fill $A_r, h_r$ \\
\ref{sec:qx} (incl.\ \ref{sec:approximations},\ref{sec:moments}) & transition at fixed $q(f)$ &
  per-transition $(\partial/\partial m_i,\partial/\partial S_i,\partial/\partial S_{i,i+1})$
  via Stein + Taylor at $m_i$ + gmix-gather drift moments \\
\ref{sec:algorithm} & — & end-to-end pseudocode + complexity \\
\hline
\end{tabular}
\end{center}

% ============================================================================
\section{The SING contract: exactly what EFGP must supply}
\label{sec:interface}
% ============================================================================

SING's backbone (algorithm and derivation: the SING paper) asks a drift
model for \emph{three} per-transition moments and nothing else. They are
the methods of \code{sing.sde.SDE}, called at the current
$q(x_i)=\mathcal{N}(m_i,S_i)$:

\begin{center}
\renewcommand{\arraystretch}{1.2}
\begin{tabular}{llll}
\hline
\code{SDE} method & math & shape & name in \code{trm} \\\hline
\code{f(\ldots, t, m, S)}    & $\Eqx[f(x_i)]$          & $(D)$   & \code{Ef}    \\
\code{ff(\ldots, t, m, S)}   & $\Eqx[\|f(x_i)\|^2]$    & scalar  & \code{Eff}   \\
\code{dfdx(\ldots, t, m, S)} & $\Eqx[J_f(x_i)]$        & $(D,D)$ & \code{Edfdx} \\\hline
\end{tabular}
\end{center}

\noindent
The three land in \code{compute\_neg\_CE\_single}
(\code{sing/utils/sing\_helpers.py:289--304}), which assembles the
per-transition negative cross-entropy into a scalar \code{trm}
(eq.~\ref{eq:transition-frozen} and \S\ref{sec:qx}).
\code{nat\_grad\_transition} (\code{sing/sing.py:46}) then
\code{vmap(jax.grad(\ldots))}s that scalar w.r.t.\ the \emph{mean}
parameters $\mu=(\E x_i,\, \E x_i x_i^\top,\, \E x_i x_{i+1}^\top)$;
by exponential-family duality that \emph{is} the natural gradient on
$(J,h,L)$, which SING $\rho$-damps and applies.

\paragraph{Two conventions to keep straight.}
(i) The diffusion enters \code{trm} as one scalar prefactor
$-1/(2\Delta_i\sigma^2)$ (line 304): the code is the \textbf{isotropic}
case $\Sigma=\sigma^2 I$ (\code{sigma\_drift\_sq} is a scalar), so the
$\Sigma^{-1}$ / $\sigma_r^{-2}$ written throughout this note is a single
$\sigma^{-2}$ in practice, and \code{ff} carries no $\Sigma^{-1}$ inside.
(ii) Lines 301--302 add the linear-input contribution $Bu_i$; the second
of these also multiplies \code{Ef}, so \code{Ef}'s cotangent picks up a
$+2\Delta_i^2 (Bu_i)$ term. Neither changes any statement below.

\paragraph{The contract is three values plus six VJP slots.}
Since SING only ever differentiates \emph{through} the three returns, a
drift backend is fully specified per transition by those three values
and their VJPs w.r.t.\ $(m_i,S_i)$. With
$\bar f_r(x)=\phi(x)^*\mu_r$ the $q(f)$ posterior-mean drift and $g$ the
upstream cotangent, EFGP supplies:

\begin{center}
\renewcommand{\arraystretch}{1.35}
\footnotesize
\begin{tabular}{@{}p{1.4cm}|p{3.5cm}|p{3.3cm}|p{4.1cm}|p{2.5cm}@{}}
\hline
slot & value (how) & $\partial/\partial m_i$ & $\partial/\partial S_i$ & status \\\hline
\code{Ef}
 & $\Eqx[\bar f]$ (gmix-gather)
 & $\Eqx[J_{\bar f}]^\top g$
 & $0$
 & $m$ exact; $S$ GN-dropped \\
\code{Eff}
 & $\|\Eqx[\bar f]\|^2$ (from \code{Ef})
 & $2g\,\Eqx[J_{\bar f}]^\top\Eqx[\bar f]$
 & $g\,\Eqx[J_{\bar f}]^\top\Eqx[J_{\bar f}]\;\succeq 0$
 & GN $+$ cov-drop \\
\code{Edfdx}
 & $\Eqx[J_{\bar f}]$ (gmix-gather)
 & $0$ (raw array)
 & $0$ (raw array)
 & GN-dropped \\\hline
\end{tabular}
\end{center}

\noindent
Code: \code{FrozenEFGPDrift} with \code{ef\_with\_jac\_grad} and
\code{eff\_with\_grads} (\code{sing/efgp\_jax\_drift.py:39--134}).
Note \code{Eff}'s forward stores only $\|\Eqx[\bar f]\|^2$; the
$\tr[\,\cdot\,S_i]$ piece of the true value lives only in the
$\partial/\partial S_i$ slot, which is all SING consumes.

\paragraph{Departures from the generic SING algorithm.}
\begin{itemize}\setlength\itemsep{3pt}
\item[\textbf{D1}] \textbf{$q(f)$ is a Fourier-feature Gaussian updated in
closed form.} Generic SING takes a fixed parametric drift; EFGP's drift
is a GP whose weight posterior $q(w_r)=\mathcal{N}(\mu_r, A_r^{-1})$
solves $A_r\mu_r=h_r$ --- an extra coordinate-ascent block interleaved
with SING's $q(x)$ update (\S\ref{sec:qf-update}--\ref{sec:compute}), plus
a collapsed kernel M-step (\texttt{efgp\_mstep.tex}).
\item[\textbf{D2}] \textbf{The $S$-slots are Gauss--Newton truncated}
(\S\ref{sec:qx}). Generic SING would let \code{jax.grad} carry the exact
Price/Bonnet derivatives through an on-graph quadrature (e.g.\
\code{GaussHermiteQuadrature}). EFGP computes the moments off-graph
(autodiff through NUFFTs inside the per-transition \code{vmap} is
${\sim}50\times$; App.~\ref{app:impl}) and injects the GN-truncated table
above --- which is what buys a PSD-by-construction $\partial/\partial S_i$
on the load-bearing slot.
\item[\textbf{D3}] \textbf{\code{ff} carries only $\bar f$, not the $q(f)$
variance.} The exact $q(f)$-average is
$\Eq[q(f)][\|f(x)\|^2]=\|\bar f(x)\|^2 + V(x)$; production drops $V$ and
its gradient (\code{restore\_qf\_variance='none'}, ``Approximation A'',
\S\ref{sec:qx} item (iv)). Opt-in restoration: App.~\ref{app:vhutch}.
\end{itemize}
Everything else in the ELBO --- emissions, $H[q(x)]$, initial --- is
Gaussian closed-form and handled unchanged by SING's backbone. A reader
who wants only the interface can stop here.

% ============================================================================
\section{Notation cheat-sheet}
\label{sec:notation}
% ============================================================================

\noindent
\begin{tabular}{@{}p{0.30\textwidth} p{0.66\textwidth}@{}}
$m_i, S_i$ & marginal mean and covariance of $q(x_i)$. \\
$S_{i,i+1}$ & cross-covariance $\Cov(x_i, x_{i+1})$. \\
$d_i := m_{i+1}-m_i$ & mean increment. \\
$C_i := S_{i,i+1}-S_i = \Cov(x_i, x_{i+1}-x_i)$ & Stein cross-cov correction. \\
$\phi(x), \phi_\theta(x)$ & finite basis for $f$; subscript $\theta$ added
  when the kernel parameterisation matters (§\ref{sec:efgp} onward). \\
$\xi_k, D_\theta$ & equispaced Fourier spectral grid and diagonal
  kernel-weight matrix (eq.~\ref{eq:efgp-features}). \\
$A_r, h_r, \mu_r$ & $q(f_r)$ precision, RHS, and posterior mean in the
  basis (eq.~\ref{eq:Ah-system}). \\
$\bar f_r(x) = \phi(x)^*\mu_r$ & posterior mean drift under $q(f)$
  (exact). \\
$\widetilde f(x_i)$ & first-order Taylor of $\bar f$ at $m_i$, the
  §\ref{sec:approximations} linearisation. \\
$J_{\bar f}(x)$ & Jacobian of $\bar f$ at $x$; for Fourier features
  cheap (eq.~\ref{eq:phi-jac}). \\
$V_r(x) = \sigma_r^{-2}\phi(x)^* A_r^{-1}\phi(x)$ & $q(f)$ posterior
  variance contribution at $x$ (eq.~\ref{eq:Eff-full}). \\
$\omega_r(\delta)$ & BTTB-diagonal projection of $D_\theta A_r^{-1}
  D_\theta$ on lag $\delta$ (App.~\ref{app:vhutch}, eq.~\ref{eq:omega-def}). \\
\end{tabular}

% ============================================================================
\section{The q(f) update: GP regression on latent velocities}
\label{sec:qf-update}
% ============================================================================

\elbobanner{transition term $\sum_i \E_{q(x,f)}\!\bigl[\log p(x_{i+1}\given x_i,f)\bigr] -
\mathrm{KL}[q(f)\|p(f)]$ at fixed $q(x)$.}{closed-form Gaussian
$q(w_r)=\mathcal{N}(\mu_r, A_r^{-1})$ via $A_r\mu_r = h_r$
(eq.~\ref{eq:Ah-system}).}

% ----------------------------------------------------------------------------
\subsection{Transition stripped to its $f$-content}
\label{sec:transition}
% ----------------------------------------------------------------------------

The Euler transition log-density is
\begin{equation}
\log p(x_{i+1}\given x_i,f)
\;=\;
-\,\frac{1}{2\Delta_i}
\bigl[\delta_i - \Delta_i\, f(x_i)\bigr]^{\!\top}\!
\Sigma^{-1}\,
\bigl[\delta_i - \Delta_i\, f(x_i)\bigr]
\;+\; \text{const}.
\end{equation}
The constant and the $\delta_i^\top\Sigma^{-1}\delta_i$ term do not involve $f$.
Expanding for diagonal $\Sigma$, the part of $\log p(x_{i+1}\given x_i,f)$
that depends on coordinate $f_r$ is
\begin{equation}
\boxed{\;
\frac{1}{\sigma_r^2}\,\delta_{i,r}\, f_r(x_i)
\;-\;
\frac{\Delta_i}{2\sigma_r^2}\, f_r(x_i)^2.
\;}
\label{eq:f-quadratic}
\end{equation}
This is a Gaussian likelihood, quadratic in the function values $f_r(x_i)$.
That is the \emph{only} fact about the transition we will need. The $q(f_r)$
update reduces to combining \eqref{eq:f-quadratic} with the GP prior
\eqref{eq:gp-prior}. Because $\Sigma$ is diagonal and the $f_r$-priors are
independent, the $q(f_r)$ updates decouple across $r$; from now on we work
coordinate-wise and suppress the index $r$ when convenient.

% ----------------------------------------------------------------------------
\subsection{Layer 1: GP regression with known $x$}
\label{sec:known-x}
% ----------------------------------------------------------------------------

Pretend, for this section only, that the latent path
$X = \{x_0,\ldots,x_{T-1}\}$ is observed. Define the pseudo-velocity
\begin{equation}
y_{i,r}^{\text{pseudo}}
\;:=\; \frac{\delta_{i,r}}{\Delta_i}
\;=\; \frac{x_{i+1,r}-x_{i,r}}{\Delta_i}.
\end{equation}
Completing the square in \eqref{eq:f-quadratic} shows the contribution from
transition $i$ is
\begin{equation}
-\,\frac{\Delta_i}{2\sigma_r^2}
\bigl(y_{i,r}^{\text{pseudo}} - f_r(x_i)\bigr)^{\!2}
\;+\; \text{const},
\end{equation}
i.e.\ a Gaussian observation of $f_r(x_i)$ with precision
$\lambda_{i,r} = \Delta_i/\sigma_r^2$. Stack
\begin{equation}
f_{r,X} = \bigl[f_r(x_0),\ldots,f_r(x_{T-1})\bigr]^{\!\top}\!,
\qquad [K_X]_{ij} = k_\theta(x_i,x_j),
\qquad \Lambda_r = \diag(\lambda_{i,r}),
\end{equation}
and let $y_r^{\text{pseudo}}$ collect the pseudo-velocities. Then
$q(f_{r,X}) = \mathcal{N}(m_{r,X},\,C_{r,X})$ with the standard GP-regression
identities
\begin{equation}
\boxed{\;
C_{r,X}^{-1} \;=\; K_X^{-1} + \Lambda_r,
\qquad
C_{r,X}^{-1}\, m_{r,X} \;=\; \Lambda_r\, y_r^{\text{pseudo}}.
\;}
\label{eq:known-x}
\end{equation}

\medskip
\noindent
\textbf{Template.} \emph{The $q(f)$ update is GP regression on latent
velocities.} The remainder of this section modifies \eqref{eq:known-x} to
the case where $x$ is uncertain; §\ref{sec:efgp} then specialises the basis
to make it FFT-cheap.

% ----------------------------------------------------------------------------
\subsection{Layer 2: uncertain $x$, in a finite basis}
\label{sec:uncertain-x}
% ----------------------------------------------------------------------------

In SING $X$ is not observed; only $q(x_{0:T})$ is available, a Gaussian
Markov variational posterior with marginals $(m_i, S_i)$ and
cross-covariances $S_{i,i+1}$. The K-space packaging \eqref{eq:known-x}
no longer applies directly --- with $X$ random, both $f_{r,X}$ and $K_X$
are random --- so rather than carry a function posterior we parametrise
the drift directly by weights,
\begin{equation}
f_r(x) \;=\; \phi(x)^{*}\, w_r,
\qquad w_r \sim \mathcal{N}(0, I_M),
\label{eq:basis}
\end{equation}
the basis chosen so the iid prior on $w_r$ reproduces the GP prior on
$f_r$. The coordinate-ascent VI update for $q(w_r)$ is then \emph{just
the expected log-joint as a function of $w_r$}:
\begin{equation}
\log q^*(w_r)
\;=\;
\Eq[q(x)]\!\bigl[\log p(x_{0:T},\, w_r)\bigr]
\;+\; \text{const},
\label{eq:logqstar}
\end{equation}
the expectation over $x$ only, with $w_r$ the variable being updated.
The only $w_r$-dependent pieces are the prior
$\log p(w_r) = -\tfrac12\, w_r^{*}w_r$ and the transition content
\eqref{eq:f-quadratic} with $f_r(x_i)=\phi(x_i)^{*}w_r$ --- so
$\delta_{i,r}\,f_r(x_i)=\delta_{i,r}\,\phi(x_i)^{*}w_r$ is linear in
$w_r$ and $f_r(x_i)^2=w_r^{*}\phi(x_i)\phi(x_i)^{*}w_r$ is quadratic.
Hence
\begin{equation}
\log q^*(w_r)
\;=\;
-\tfrac12\, w_r^{*} A_r\, w_r
\;+\;
\Re\!\bigl(h_r^{*} w_r\bigr)
\;+\;\text{const},
\end{equation}
with
\begin{equation}
\boxed{\;
A_r \;=\; I + \frac{1}{\sigma_r^2}\sum_i \Delta_i\,
\Eqx\!\bigl[\phi(x_i)\phi(x_i)^{*}\bigr],
\qquad
h_r \;=\; \frac{1}{\sigma_r^2}\sum_i
\Eqxx\!\bigl[\phi(x_i)\,(x_{i+1,r}-x_{i,r})\bigr].
\;}
\label{eq:Ah}
\end{equation}
The expectation over $q(x)$ touches only the features, contracting
$\phi\phi^{*}$ and $\phi$ to the deterministic $M\times M$ matrix and
$M$-vector in \eqref{eq:Ah}; $w_r$ is never integrated. So
$\log q^*(w_r)$ is exactly quadratic in $w_r$, and completing the square
gives the Gaussian
\begin{equation}
q(w_r) \;=\; \mathcal{N}(\mu_r,\,A_r^{-1}),
\qquad
A_r\, \mu_r \;=\; h_r.
\label{eq:Ah-system}
\end{equation}

\paragraph{Stein collapse of the linear term in $h_r$.}
The pairwise expectation in \eqref{eq:Ah} under $q(x_i, x_{i+1})$ is
awkward. Stein's lemma collapses it to marginal expectations plus a
cross-cov correction. With $d_i := m_{i+1}-m_i$ and $C_i := S_{i,i+1}-S_i
= \Cov(x_i,\, x_{i+1}-x_i)$, for any smooth $g$,
\begin{equation}
\E\bigl[g(x_i)\,(x_{i+1,r}-x_{i,r})\bigr]
\;=\;
\E[g(x_i)]\,d_{i,r}
\;+\;
\E[J_g(x_i)]\,C_{i,\cdot r}.
\end{equation}
Apply with $g=\phi$ (componentwise) inside \eqref{eq:Ah}:
\begin{equation}
\boxed{\;
h_r
\;=\;
\frac{1}{\sigma_r^2}\sum_i
\Bigl\{
\underbrace{\Eqx[\phi(x_i)]\,d_{i,r}}_{\text{mean-increment}}
\;+\;
\underbrace{\Eqx[J_\phi(x_i)]^{*}\,C_{i,\cdot r}}_{\text{Stein correction}}
\Bigr\}.
\;}
\label{eq:hr-stein-abstract}
\end{equation}
Every quantity in \eqref{eq:Ah-system} now depends only on \emph{marginal}
expectations under $q(x_i)$, plus the cross-cov correction $C_i$.

% ============================================================================
\section{The EFGP basis choice}
\label{sec:efgp}
% ============================================================================

\elbobanner{still the transition term at fixed $q(x)$.}{Fourier basis
specialisation of \eqref{eq:Ah-system}: $A_r$ becomes BTTB-Toeplitz and
$J_\phi$ is frequency multiplication.}

EFGP specialises the basis $\phi$ in \eqref{eq:basis} to Fourier features:
\begin{equation}
\phi_\theta(x) \;=\; D_\theta\, F_x,
\qquad
F_x[k] \;=\; e^{2\pi i\, x^\top \xi_k},
\qquad
\phi_\theta(x)^{*}\phi_\theta(x') \;\approx\; k_\theta(x,x'),
\label{eq:efgp-features}
\end{equation}
with $\{\xi_k\}_{k=1}^{M}$ an equispaced spectral grid and $D_\theta$
diagonal with square-root spectral weights chosen to match $k_\theta$.
This is the only kernel approximation introduced. The system
\eqref{eq:Ah-system} is unchanged in form; what we gain are two
structural simplifications.

\paragraph{Cheap Jacobian.}
Differentiating \eqref{eq:efgp-features},
\begin{equation}
\frac{\partial \phi_k(x)}{\partial x_j} \;=\; 2\pi i\,\xi_{k,j}\,\phi_k(x),
\label{eq:phi-jac}
\end{equation}
so $J_\phi$ is just frequency-multiplication of $\phi$. The Stein
correction in \eqref{eq:hr-stein-abstract} costs no more than evaluating
$\phi$ itself.

\paragraph{BTTB structure of $\Eqx[\phi\phi^{*}]$.}
For Fourier features, $\phi(x)\phi(x)^{*}$ has $(k,\ell)$ entry
$D_{\theta,k}\,e^{2\pi i x^\top(\xi_\ell-\xi_k)}\,D_{\theta,\ell}$. The
expectation under $q(x_i)$ depends on $\xi_\ell-\xi_k$ only:
\begin{equation}
\bigl(\Eqx[\phi\phi^{*}]\bigr)_{k\ell}
\;=\;
D_{\theta,k}\,\Eqx\!\bigl[e^{2\pi i\,x_i^\top(\xi_\ell-\xi_k)}\bigr]\,D_{\theta,\ell}.
\end{equation}
On an equispaced grid this makes $\Eqx[\phi\phi^{*}]$ Block-Toeplitz with
Toeplitz Blocks (BTTB), so $A_r v$ costs $O(M\log M)$ via FFT instead of
$O(M^2)$.

% ============================================================================
\section{Computing the $q(x)$-expectations that build $A_r, h_r$}
\label{sec:compute}
% ============================================================================

\elbobanner{still the transition term at fixed $q(x)$.}{the marginal
expectations $\Eqx[\phi(x_i)\phi(x_i)^*]$, $\Eqx[\phi(x_i)]$,
$\Eqx[J_\phi(x_i)]$ that fill \eqref{eq:Ah-system}–\eqref{eq:hr-stein-abstract}.}

Because every $q(x_i)=\mathcal{N}(m_i,S_i)$ is Gaussian and every $\phi_k$
is a complex exponential, each marginal expectation appearing in
\eqref{eq:Ah}--\eqref{eq:hr-stein-abstract} is the Gaussian characteristic
function evaluated at frequency differences:
\begin{equation}
\Eqx\!\bigl[e^{2\pi i\,x_i^\top\xi}\bigr]
\;=\;
e^{2\pi i\,m_i^\top\xi}\;e^{-2\pi^2\,\xi^\top S_i\,\xi},
\label{eq:Eqx-chi}
\end{equation}
a phase at $m_i$ times a Gaussian envelope set by $S_i$. So \emph{every}
quantity needed in \eqref{eq:Ah-system}--\eqref{eq:hr-stein-abstract} is
one \textbf{Gaussian-enveloped non-uniform Fourier sum}
\begin{equation}
F(\xi)\;=\;\sum_i w_i\, e^{2\pi i\, m_i^\top \xi}\; e^{-2\pi^2\,\xi^\top S_i\,\xi},
\label{eq:enveloped-nufft}
\end{equation}
evaluated on the frequency-\emph{difference} grid (for the BTTB generator
of $A_r$) or on the spectral grid (for $h_r$), with $w_i$ the relevant
Stein weight ($\Delta_i/\sigma^2$, $d_{i,r}/\sigma^2$, or
$[C_i]_{j,r}/\sigma^2$; assembled by \code{\_flatten\_stein}). Three
interchangeable backends compute \eqref{eq:enveloped-nufft}, differing
only in how the \emph{per-source} envelope $S_i$ is handled.
\code{estep\_method='auto'} (the default) resolves by device:

\begin{center}
\renewcommand{\arraystretch}{1.3}
\footnotesize
\begin{tabular}{@{}p{2.2cm}|p{4.9cm}|p{3.1cm}|p{4.4cm}@{}}
\hline
\code{estep\_method} & envelope handling & cost & error \\\hline
\code{'gmix'} \emph{(default on GPU)}
 & each $\mathcal{N}(m_i,S_i)$ tabulated on a fine
   spatial grid over a $(2r{+}1)^D$ stencil, then FFT
 & $O(N\,r^D + N_f^D\log N_f)$
 & stencil tail truncation ($n_\sigma$) \\
\code{'analytic'} \emph{(default on CPU; $D{=}2$ only)}
 & $S_i = \bar S + \Delta S_i$; Taylor
   $e^{-2\pi^2\xi^\top\Delta S_i\xi}$ to order $p$
 & $O(P(p)\,(N + M\log M))$
 & $\Delta S$ truncation; \emph{exact} at $\Delta S{=}0$ \\
\code{'mc'} \emph{(reference)}
 & $S$ marginal samples per source, plain NUFFT
 & $O(NS + M\log M)$
 & Monte Carlo variance \\\hline
\end{tabular}
\end{center}

\noindent
Resolution logic (\code{efgp\_em.py:698--712}): \code{'gpu'}$\to$\code{'gmix'},
otherwise \code{'analytic'}, falling back to \code{'gmix'} when $D\ne 2$.
The devices disagree because the bottleneck does: on CPU the gmix
per-source spread dominates the wall, so \code{'analytic'} is ${\sim}6\times$
faster on the $q(f)$ update; on GPU scatter and cuFFT are nearly free
(the wall is JIT compile plus the sequential-in-$T$ block-tridiagonal
smoother \code{lax.scan}), so \code{'analytic'} only adds cuFINUFFT
launch overhead. Any explicit value overrides \code{'auto'}.

% ----------------------------------------------------------------------------
\subsection{Backend 1: gmix closed-form spreader (GPU default)}
\label{sec:gmix}
% ----------------------------------------------------------------------------

Substituting
into $A_r$ \eqref{eq:Ah} preserves BTTB structure: the generating kernel
on the difference grid $\Delta\xi_k$ is the Type-1 NUFFT of weighted
Gaussians \emph{centred at $m_i$ with per-source widths $S_i$}.
Operationally this replaces the fixed NUFFT spreading kernel by a
per-source $\mathcal{N}(\,\cdot\,;m_i,S_i)$ on a truncated stencil;
the FFT and deconvolve steps are unchanged. The RHS terms in
\eqref{eq:hr-stein-abstract} receive the same treatment.

The companion note \texttt{efgp\_estep\_charfun.tex} carries the full
derivation, the explicit closed-form generator, complexity proofs, and
the bucketed / Gauss--Hermite alternatives considered. The closed-form
spreader eliminates Monte Carlo variance at the cost of a per-source
spread; in single-scale regimes the asymptotic cost matches a standard
NUFFT, and heterogeneous $S_i$ costs an extra
$\propto(\sigma_{\max}/\sigma_{\min})^{D}$ factor. The stencil is
truncated at $n_\sigma$ standard deviations (default $1.5$), which leaves
a tail bias present \emph{even at} $\Delta S=0$ (${\sim}10\%$ vs.\ MC on
the reference problem) --- the accuracy motivation for Backend 2.
Code: \code{compute\_mu\_r\_gmix\_jax}
(\code{efgp\_jax\_drift.py:269}), spreader in
\code{sing/efgp\_gmix\_spreader.py}.

% ----------------------------------------------------------------------------
\subsection{Backend 2: analytic Taylor-envelope NUFFT (CPU default)}
\label{sec:analytic}
% ----------------------------------------------------------------------------

Split the per-source covariance about the trajectory mean,
$S_i = \bar S + \Delta S_i$, and Taylor-expand \emph{only} the deviation
factor:
\begin{equation}
e^{-2\pi^2\,\xi^\top S_i\,\xi}
= e^{-2\pi^2\,\xi^\top \bar S\,\xi}\, e^{u_i(\xi)},
\quad u_i(\xi):=-2\pi^2\,\xi^\top \Delta S_i\,\xi,
\quad e^{u_i}\approx\sum_{k\le p}\frac{u_i^k}{k!}.
\label{eq:env-taylor}
\end{equation}
Each term of $u_i^k$ factorises into a $\xi$-monomial times a
$\Delta S_i$-monomial, so \eqref{eq:enveloped-nufft} becomes a fixed
number of \emph{ordinary} type-1 NUFFTs on the same nodes $\{m_i\}$ with
re-weighted source vectors:
\begin{equation}
F(\xi)\;\approx\;e^{-2\pi^2\xi^\top\bar S\xi}
\sum_{\rm pairs} c_{\rm pair}(\xi)\;
\mathrm{NUFFT1}\bigl(\{m_i\},\,\{w_i\,\nu_{{\rm pair},i}\}\bigr)(\xi),
\label{eq:analytic-final}
\end{equation}
with the $(c_{\rm pair},\nu_{\rm pair})$ enumerated by
\code{\_env\_taylor\_pairs} (\code{efgp\_jax\_drift.py:525}):
$P(p)=\binom{p+2}{2}$ pairs in 2-D, i.e.\ $1, 4, 10$ at
\code{analytic\_order} $p=0,1,2$. All pairs share the nodes $\{m_i\}$ and
are evaluated as ONE batched transform (a single FINUFFT plan), so
order$\to$order cost is nearly flat; the same expansion is applied on the
difference grid (BTTB generator) and on the spectral grid (RHS).
Code: \code{compute\_mu\_r\_analytic\_jax}
(\code{efgp\_jax\_drift.py:569}).

No fine grid and no stencil: \emph{exact} for homogeneous $S$ (any order),
with the truncation error controlled by the $S$-spread rather than by a
grid. Order 1 holds to ${<}4\%$ up to ${\sim}20\times$ spread in the
$S_i$ eigenvalues, order 2 to ${\sim}60\times$; converged fits sit near
$1.4\times$. 2-D only in v0. Full derivation and benchmarks:
\texttt{efgp\_analytic\_qf.tex}.

% ----------------------------------------------------------------------------
\subsection{Backend 3: Monte Carlo cloud (reference)}
\label{sec:mc}
% ----------------------------------------------------------------------------

The same expectations can be evaluated by Monte Carlo. For each
$i=0,\ldots,T-2$, draw $S$ marginal samples $x_i^{(s)}\sim q(x_i)$ and
stack as $X=\{x_i^{(s)}\}_{i,s}\in\mathbb{R}^{(T-1)S\times D_x}$. With
\begin{equation}
[F_X]_{k,(i,s)} \;=\; e^{2\pi i\,(x_i^{(s)})^\top \xi_k},
\qquad
\phi_\theta(x_i^{(s)}) \;=\; D_\theta\,(F_X)_{:,(i,s)},
\end{equation}
and per-sample weights $(W_r)_{(i,s)} = \Delta_i/(S\sigma_r^2)$,
$A_r$ becomes
\begin{equation}
A_r \;\approx\; I + D_\theta\,T_r\,D_\theta,
\qquad
T_r \;:=\; F_X^{*}\,W_r\,F_X,
\end{equation}
with $(T_r)_{k\ell}$ depending on $\xi_\ell-\xi_k$ only. $T_r$ is BTTB on
the equispaced grid; a single Type-1 NUFFT over $X$ produces its generating
kernel, and $A_r v = v + D_\theta\,(T_r\,(D_\theta v))$. The system
$A_r\mu_r=h_r$ is solved by Jacobi-preconditioned CG.

For $h_r$, set
\begin{equation}
(a_r)_{(i,s)} = \frac{d_{i,r}}{S\,\sigma_r^2},
\qquad
(c_{j,r})_{(i,s)} = \frac{[C_i]_{j,r}}{S\,\sigma_r^2},
\end{equation}
so \eqref{eq:hr-stein-abstract} estimates as
\begin{equation}
h_{1,r} \approx D_\theta\, F_X^{*} a_r,
\qquad
h_{2,r} \approx \sum_{j=1}^{D_x} (-2\pi i\,\xi_j) \odot D_\theta\, F_X^{*} c_{j,r},
\qquad
h_r = h_{1,r} + h_{2,r}.
\label{eq:hr-mc}
\end{equation}
The factor $-2\pi i\,\xi_j$ in $h_{2,r}$ (rather than $+2\pi i\,\xi_j$)
comes from the adjoint Jacobian $J_\phi^{*}$:
$\overline{2\pi i\,\xi_{k,j}\,\phi_k(x)} = -2\pi i\,\xi_{k,j}\,\overline{\phi_k(x)}$.

The Stein form \eqref{eq:hr-stein-abstract} can be checked against a naive
pair-sampling estimator (draw transition pairs $(x_i^{(s)}, x_{i+1}^{(s)})
\sim q(x_i, x_{i+1})$ and use sampled pseudo-velocities in
\eqref{eq:known-x}); both estimators target the same $h_r$, with Stein
strictly lower variance. SING uses Stein in production and pair-sampling
only as a regression test.

% ============================================================================
\section{The q(x) natural-gradient step}
\label{sec:qx}
% ============================================================================

\elbobanner{transition term at fixed $q(f)$ ---
$\sum_i \Eqxx\bigl[\Eq[q(f)][\log p(x_{i+1}\given x_i, f)]\bigr]$.}{per-transition
$(\partial/\partial m_i,\, \partial/\partial S_i,\, \partial/\partial S_{i,i+1})$
of this term, ready for SING's $\rho$-damped natural-grad update on $q(x)$.}

\subsection*{Setup}

SING does $\rho$-damped natural gradient on $q(x)$ (constrained
Gaussian-Markov), using
\begin{equation}
\widetilde\nabla_\eta \mathcal{L} \;=\; \partial \mathcal{L}/\partial \mu,
\qquad \eta=(J,h,L),\quad
\mu=(\E[x_i], \E[x_ix_i^\top], \E[x_ix_{i+1}^\top]),
\label{eq:natgrad-duality}
\end{equation}
i.e.\ exponential-family duality (\S\ref{sec:interface}). Mean params
re-parameterise cleanly to $(m_i, S_i, S_{i,i+1})$. This section derives
\emph{why} the six VJP slots of \S\ref{sec:interface} take the values they
do --- departure \textbf{D2}.

\subsection*{The transition term, and how the code assembles it}

At fixed $q(f)$, with $\bar f_r(x) = \phi(x)^* \mu_r$ deterministic but
nonlinear in $x$ and $V_r(x) = \sigma_r^{-2}\phi(x)^* A_r^{-1}\phi(x)$
the $q(f)$ posterior variance contribution, the per-transition piece is
\begin{equation}
\begin{aligned}
\Eqxx\!\bigl[\Eq[q(f)][\log p(x_{i+1}\given x_i, f)]\bigr]
\;=\;&
-\,\frac{\E[\delta_i^\top\Sigma^{-1}\delta_i]}{2\Delta_i}
+ \Eqxx[\delta_i^\top\Sigma^{-1}\bar f(x_i)]\\
&- \frac{\Delta_i}{2}\,\Eqx\!\bigl[\bar f^\top\Sigma^{-1}\bar f + V\bigr]
+ \text{const}.
\end{aligned}
\label{eq:transition-frozen}
\end{equation}
This is assembled into the scalar \code{trm} of
\code{compute\_neg\_CE\_single} (\code{sing\_helpers.py:289--304}) from
four pieces:

\begin{center}
\renewcommand{\arraystretch}{1.3}
\begin{tabular}{l l l}
\hline
piece & math & where \code{trm} is built (lines in \code{compute\_neg\_CE\_single}) \\\hline
(i)   & $\E[\delta^\top\Sigma^{-1}\delta]/(2\Delta_i)$            & 294-296 (traces of $m, S, SS$) \\
(ii)  & $\Eqxx[\delta^\top\Sigma^{-1}\bar f]$                     & 298-299 (\code{Ef} and \code{Edfdx}) \\
(iii) & $\Eqx[\bar f^\top\Sigma^{-1}\bar f]$                      & 297 (\code{(del\_t)**2 * Eff}) \\
(iv)  & $\Eqx[V]$                                                  & dropped (production) \\\hline
\end{tabular}
\end{center}

The drift moments \code{Ef}, \code{Eff}, \code{Edfdx} consumed in
lines 297-299 are pre-computed off-graph by
\code{drift\_moments\_gmix\_jax} and exposed via the
\code{FrozenEFGPDrift} pytree --- from
\code{compute\_neg\_CE\_single}'s view they are per-transition
constants. The Taylor framework below dictates how their
$(m_t, S_t)$-gradients are filled in via \code{custom\_vjp}; the
four item subsections then walk each piece's value and gradient flow.

\subsection*{The Gauss-Newton framework}

Items (ii)-(iv) need $(m_i, S_i)$-derivatives of expectations under
$q(x_i) = \mathcal{N}(m_i, S_i)$ of nonlinear functions ($\bar f$,
$J_{\bar f}$, $\bar f^\top \Sigma^{-1}\bar f$, $V$). Production handles
all of these under a single principled approximation: a Gauss-Newton
(GN) truncation of Price's theorem.

\paragraph{Price's theorem (Bonnet's lemma).}
For any sufficiently smooth $g$ and Gaussian $q(x_i)$,
\begin{equation}
\partial_{m_i} \E_{q(x_i)}[g(x_i)] = \E_{q(x_i)}[\nabla g(x_i)],
\qquad
\partial_{S_i} \E_{q(x_i)}[g(x_i)] = \tfrac{1}{2}\E_{q(x_i)}[\nabla^2 g(x_i)].
\label{eq:price}
\end{equation}
These are exact identities, expressing the $(m, S)$-derivatives of a
Gaussian expectation as expectations of the gradient and Hessian of
the integrand.

\paragraph{The GN move on the load-bearing slot.}
For $g = \bar f^\top \Sigma^{-1}\bar f$ (the quadratic in item (iii)),
the Hessian splits as
\begin{equation}
\nabla^2 \!\bigl(\bar f^\top \Sigma^{-1}\bar f\bigr) \;=\; 2\bigl(
\underbrace{J_{\bar f}^\top \Sigma^{-1} J_{\bar f}}_{\text{GN piece: PSD}}
+ \underbrace{\textstyle\sum_r \sigma_r^{-2}\,\bar f_r\,H_{\bar f_r}}_{\text{residual: sign-indefinite}}\bigr).
\label{eq:gn-hessian}
\end{equation}
The first piece is PSD-by-construction; the second (the residual
Hessian) is sign-indefinite because $\bar f_r$ has no sign constraint.
\textbf{The Gauss-Newton approximation drops the residual Hessian} ---
the standard move for sum-of-squares objectives. Plugged into Price:
\begin{equation}
\begin{aligned}
\partial_{S_i} \E[\bar f^\top \Sigma^{-1}\bar f]
&\;\stackrel{\rm Price}{=}\; \E\bigl[J_{\bar f}^\top \Sigma^{-1} J_{\bar f}\bigr]
   + \tfrac{1}{2}\E\!\bigl[\textstyle\sum_r \sigma_r^{-2}\bar f_r H_{\bar f_r}\bigr]
\;\stackrel{\rm GN}{\approx}\;
\E\bigl[J_{\bar f}^\top \Sigma^{-1} J_{\bar f}\bigr]\\
&\;\stackrel{\rm Cov\text{-}drop}{\approx}\;
\E[J_{\bar f}]^\top \Sigma^{-1}\,\E[J_{\bar f}] \;\succeq\; 0.
\end{aligned}
\label{eq:Eff-lin}
\end{equation}
The second approximation (the covariance drop) replaces $\E[J^\top J]$
with $\E[J]^\top \E[J]$ --- also second-order in the nonlinearity of
$\bar f$. The result is PSD-by-construction, a structural consequence
of the GN form rather than a coincidence.

\paragraph{The same GN move handles every other truncation.}
For the linear pieces, Price + GN gives
\begin{equation}
\partial_{S_i} \E[\bar f]    \stackrel{\rm Price}{=} \tfrac{1}{2}\E[H_{\bar f}]
                                \stackrel{\rm GN}{\approx} 0,
\quad
\partial_{m_i} \E[J_{\bar f}] \stackrel{\rm Price}{=} \E[H_{\bar f}]
                                \stackrel{\rm GN}{\approx} 0,
\quad
\partial_{S_i} \E[J_{\bar f}] \stackrel{\rm Price}{=} \tfrac{1}{2}\E[\nabla^2 J_{\bar f}]
                                \stackrel{\rm GN}{\approx} 0.
\end{equation}
Every dropped slot is a residual Hessian (of $\bar f$ or of
$J_{\bar f}$) --- the same object GN drops in the quadratic case.
The cross-term gradient
$\partial_{m_i}\E[\bar f^\top \Sigma^{-1}\bar f] = 2\E[J^\top \Sigma^{-1}\bar f]$
gets a further covariance drop in production:
$2\E[J]^\top \Sigma^{-1}\E[\bar f]$.

\paragraph{Gradient table (GN-truncated Price's identities).}

\begin{center}
\renewcommand{\arraystretch}{1.4}
\footnotesize
\begin{tabular}{@{}p{2.5cm}|p{5.6cm}|p{3.3cm}|p{3.9cm}@{}}
\hline
quantity & value & $\partial/\partial m_i$ & $\partial/\partial S_i$ \\\hline
$\Eqx[\bar f]$
  & $\Eqx[\bar f]$
  & $\Eqx[J_{\bar f}]$
  & $0$ \emph{(GN: $\tfrac{1}{2}\E[H_{\bar f}]$ dropped)} \\
$\Eqx[J_{\bar f}]$
  & $\Eqx[J_{\bar f}]$
  & $0$ \emph{(GN: $\E[H_{\bar f}]$ dropped)}
  & $0$ \emph{(GN)} \\
$\Eqx[\bar f^\top\Sigma^{-1}\bar f]$
  & $\Eqx[\bar f]^\top\Sigma^{-1}\Eqx[\bar f] + \tr[\Eqx[J_{\bar f}]^\top\Sigma^{-1}\Eqx[J_{\bar f}]\,S_i]$
  & $2\,\Eqx[J_{\bar f}]^\top\Sigma^{-1}\Eqx[\bar f]$
  & $\Eqx[J_{\bar f}]^\top\Sigma^{-1}\Eqx[J_{\bar f}] \;\succeq 0$ \\\hline
\end{tabular}
\end{center}

\paragraph{Why this is justified.}
GN is a well-studied Hessian approximation for sum-of-squares
objectives: drop the residual Hessian, keep the PSD outer-product
piece. Exact when residuals vanish ($\bar f_r \to 0$); accurate when
$\|\bar f_r\|\,\|H_{\bar f_r}\| \ll \|J_{\bar f}\|^2$. The covariance
drops $\text{Cov}(J, J)$ and $\text{Cov}(J, \bar f)$ are of the same
order. This is the same approximation underlying variational
Gauss-Newton (Khan \& Lin, 2017) and drift-linearisation in
variational SDE inference (Opper \& Archambeau, 2009).

\paragraph{Why we can't just autodiff.}
The exact $\Eqx[\bar f^\top \Sigma^{-1}\bar f]$ closes via spectral
autocorrelation of $\mu_r D_\theta$,
\begin{equation}
\Eqx[\bar f^\top \Sigma^{-1}\bar f]
\;=\; \sum_r \sigma_r^{-2}\!\sum_{k,\ell}
     \mu_{r,k}^* D_k D_\ell \mu_{r,\ell}\,
     e^{2\pi i m_i^\top \delta_{k\ell}}\,
     e^{-2\pi^2 \delta_{k\ell}^\top S_i \delta_{k\ell}},
\quad \delta_{k\ell} := \xi_\ell - \xi_k,
\label{eq:gmix-exact-Eff}
\end{equation}
but $\partial_{S_i}$ of this is a sign-indefinite sum of
$\delta\delta^\top$ matrices (the full Hessian contribution to Price)
--- it would break the PSD property SING's $J$-update relies on. The
GN-truncated form \eqref{eq:Eff-lin} is what gives the load-bearing
$\partial_{S_i} \succeq 0$. Separately, \code{jax.grad} through the
spectral closed forms costs ${\sim}50\times$ per inner iter under
\code{vmap(grad(\ldots))} (App.~\ref{app:impl}).

\paragraph{Production base values via gmix-gather.}
The expectations $\Eqx[\bar f]$ and $\Eqx[J_{\bar f}]$ in the table
slots are computed exactly in closed form via the Gaussian char-fun
\eqref{eq:Eqx-chi}: with $\bar f_r(x) = \phi(x)^* \mu_r$,
\[
\Eqx[\bar f_r] = \sum_k \mu_{r,k}\,D_{\theta,k}\,e^{2\pi i m_i^\top \xi_k}\,e^{-2\pi^2 \xi_k^\top S_i \xi_k},
\]
analogously $\Eqx[\partial_j\bar f_r]$ with frequency modulation by
$(2\pi i\,\xi_{k,j})$. The gmix-gather primitive in
\code{drift\_moments\_gmix\_jax} (\code{efgp\_jax\_drift.py:793--870};
Type-2 analog of the §\ref{sec:gmix} spreader: one batched IFFT $+$
per-source Gaussian stencil gather) delivers
$(\code{Ef}, \code{Eff}, \code{Edfdx})$ at all $T$ base points in one
shot.

\paragraph{Implementation: GN-as-\code{custom\_vjp}.}
The GN choice maps operationally to \emph{treat $\Eqx[J_{\bar f}]$ as
constant w.r.t.\ $(m_i, S_i)$ in \code{jax.grad}'s chain rule}.
Implemented in \code{FrozenEFGPDrift} via three mechanisms:
\begin{itemize}\setlength\itemsep{2pt}
\item \code{Edfdx} ($=\Eqx[J_{\bar f}]$) exposed as a raw array,
  \emph{no} \code{custom\_vjp}; chain rule through \code{Edfdx}
  returns 0. This implements the GN drops
  $\partial_m\E[J_{\bar f}] = \E[H_{\bar f}] \to 0$ and
  $\partial_S\E[J_{\bar f}] \to 0$.
\item \code{Ef} ($=\Eqx[\bar f]$) through \code{ef\_with\_jac\_grad}'s
  \code{custom\_vjp}: backward returns $\partial \code{Ef}/\partial m
  = \Eqx[J_{\bar f}]^\top g$, $\partial \code{Ef}/\partial S = 0$. The
  $\partial_m$ injection coincides with the exact Price value
  $\partial_m\Eqx[\bar f] = \Eqx[J_{\bar f}]$ (Price reduces to a
  single gradient term here, no Hessian to drop); the $\partial_S = 0$
  drops $\frac{1}{2}\E[H_{\bar f}]$, the GN move.
\item \code{Eff} ($= \|\Eqx[\bar f]\|^2$) through
  \code{eff\_with\_grads}'s \code{custom\_vjp}: backward returns the
  GN entries directly, $\partial_m = 2g(\Eqx[\bar f]\cdot\Eqx[J_{\bar f}])$
  and $\partial_S = g\,\Eqx[J_{\bar f}]^\top\Eqx[J_{\bar f}]$. The
  forward stores only $\|\Eqx[\bar f]\|^2$; the $\tr$-piece is
  absorbed into the gradient (saves a per-source $D\times D$
  contraction, since SING uses the gradient not the value).
\end{itemize}

\paragraph{Strict-Taylor as a special case.}
Replacing $\bar f$ by its first-order Taylor expansion at $m_i$,
\begin{equation}
\widetilde f(x_i) := \bar f(m_i) + J_{\bar f}(m_i)(x_i - m_i),
\label{eq:linearisation}
\end{equation}
with $J_{\bar f}(m_i)$ treated as constant w.r.t.\ $(m_i, S_i)$, gives
the same gradient table by a different route. The two derivations are
equivalent; the GN derivation is cleaner because (a) it doesn't
introduce a surrogate, (b) the PSD property of the load-bearing slot
is structural, and (c) it connects to standard variational-GN
literature.

\subsection*{Item (i): the $\delta$-quadratic --- exact}

\textbf{Math.} $\E[\delta^\top\Sigma^{-1}\delta]/(2\Delta_i)$, a Gaussian
quadratic in $\delta = x_{i+1}-x_i$ under $q(x_i, x_{i+1})$ ---
closed-form in $(m_i, m_{i+1}, S_i, S_{i+1}, S_{i,i+1})$.

\textbf{Value (code, lines 294-296).}
\begin{quote}
\code{trm  = trace(St\_next + outer(mt\_next, mt\_next))} \\
\code{trm += trace(St + outer(mt, mt))} \\
\code{trm += -2 * trace(SS + outer(mt\_next, mt))}
\end{quote}

\textbf{Gradients.} Pure \code{jax.grad} autodiff. No
\code{custom\_vjp}. SING-backbone, not EFGP-specific.

\subsection*{Item (ii): the Stein cross-term --- value exact; gradients mix exact and GN-truncated pieces}
\label{sec:moments}

\textbf{Math.} The Stein identity from §\ref{sec:uncertain-x}, applied
with $g = \bar f$, collapses the pairwise expectation \emph{exactly}
into a mean-increment piece and a Jacobian (Stein-correction) piece:
\begin{equation}
\Eqxx[\delta_i^\top\Sigma^{-1}\bar f(x_i)]
\;=\;
\underbrace{d_i^\top\Sigma^{-1}\Eqx[\bar f]}_{\text{Part A: mean-increment}}
+ \underbrace{\tr\!\bigl[\Sigma^{-1}\Eqx[J_{\bar f}]\,C_i^\top\bigr]}_{\text{Part B: Jacobian}},
\label{eq:Stein-cross}
\end{equation}
with $d_i := m_{i+1}-m_i$ and $C_i := S_{i,i+1}-S_i$ as before. The
marginals $\Eqx[\bar f]$, $\Eqx[J_{\bar f}]$ are computed by
gmix-gather in closed form (framework above) and stored as \code{Ef},
\code{Edfdx} on \code{FrozenEFGPDrift}.

\textbf{Value (code, lines 298-299).}
\begin{quote}
\code{trm += -2 * del\_t * trace(outer(Ef, mt\_next) + Edfdx @ SS.T)}\\
\code{trm +=  2 * del\_t * trace(outer(Ef, mt)      + Edfdx @ St)}
\end{quote}
The outer-product terms are Part A in code form (combining to
$-2\Delta_i\,\Eqx[\bar f]^\top d_i$); the \code{Edfdx @ ...} terms are
Part B ($-2\Delta_i\,\tr(\Eqx[J_{\bar f}] C_i^\top)$). The values
themselves are exact.

\textbf{Gradients.} \code{jax.grad} of \code{trm} assembles three
contributions per slot: \emph{direct} (autodiff through the outer
products and traces of lines 298--299, holding \code{Ef}/\code{Edfdx}
constant --- exact); \emph{Ef-chain} (\code{ef\_with\_jac\_grad}'s
injections $\partial_m\code{Ef}=\Eqx[J_{\bar f}]^\top g$, exact by Price,
and $\partial_S\code{Ef}=0$, the GN drop of
$\tfrac12\E[H_{\bar f}]$ --- exactly this would be the non-zero
envelope-derivative
$-2\pi^2\sum_k \mu_{r,k}D_k\,\xi_k\xi_k^\top e^{2\pi i m^\top\xi_k}
e^{-2\pi^2\xi_k^\top S\xi_k}$); and \emph{Edfdx-chain} (a raw constant
array, so all its slots are $0$ --- the ``$\Eqx[J]$ held constant''
rule). All values and gradients use the \emph{Gaussian-averaged}
$\Eqx[\bar f]$, $\Eqx[J_{\bar f}]$ from gmix-gather, not point-evals at
$m_i$. Resolving per slot, each contribution labelled exact /
GN-truncated:

\begin{center}
\renewcommand{\arraystretch}{1.25}
\footnotesize
\begin{tabular}{@{}p{1.8cm}|p{5.2cm}|p{4.3cm}|p{4.0cm}@{}}
\hline
slot & direct & Ef-chain & Edfdx-chain \\\hline
$\partial m_t$       & $2\Delta_i\,\Eqx[\bar f]$ \emph{(exact)} & $-2\Delta_i\,\Eqx[J_{\bar f}]^\top d_i$ \emph{(exact)} & $0$ \emph{(GN: $\E[H_{\bar f}]$ dropped)} \\
$\partial m_{t+1}$   & $-2\Delta_i\,\Eqx[\bar f]$ \emph{(exact)} & $0$ (no $m_{t+1}$ in Ef) & $0$ (Edfdx indep.\ of $m_{t+1}$; exact) \\
$\partial S_t$       & $2\Delta_i\,\Eqx[J_{\bar f}]^\top$ \emph{(exact, from $\tr(\text{Edfdx}\,S_t)$)} & $0$ \emph{(GN: env-deriv dropped)} & $0$ \emph{(GN: env-deriv dropped)} \\
$\partial S_{i,i+1}$ & $-2\Delta_i\,\Eqx[J_{\bar f}]^\top$ \emph{(exact)} & $0$ (no $S_{i,i+1}$) & $0$ (Edfdx indep.\ of $S_{i,i+1}$; exact) \\\hline
\end{tabular}
\end{center}

\textbf{Net contributions of (ii) to $\partial \code{trm}$:}
\begin{align*}
\partial m_t        &= 2\Delta_i\bigl(\Eqx[\bar f] - \Eqx[J_{\bar f}]^\top d_i\bigr), &
\partial m_{t+1}    &= -2\Delta_i\,\Eqx[\bar f], \\
\partial S_t        &= 2\Delta_i\,\Eqx[J_{\bar f}]^\top, &
\partial S_{i,i+1}  &= -2\Delta_i\,\Eqx[J_{\bar f}]^\top.
\end{align*}
\textbf{Status per slot:} $\partial m_{t+1}$ and $\partial S_{i,i+1}$
are \emph{fully exact}; $\partial m_t$ has Part A's contribution exact
but Part B's GN-truncated ($\E[H_{\bar f}]$ dropped); $\partial S_t$
has the direct piece exact but both chain contributions GN-truncated.
Linear in $\bar f$; no PSD constraint applies (the PSD slot is item
(iii)'s $\partial S_t$).

\paragraph{EFGP vs.\ SING.}
This is departure \textbf{D2} (\S\ref{sec:interface}) in detail: SING's
\code{nat\_grad\_transition} accepts whatever the backend's
$(\code{f},\code{ff},\code{dfdx})$ return, so an on-graph
Gauss--Hermite backend would give \code{jax.grad} the exact
$\partial_S\Eqx[\bar f]$; EFGP goes off-graph to avoid autodiff through
NUFFTs (App.~\ref{app:impl}) and pays for it with the GN drops above.

\subsection*{Item (iii): the marginal quadratic --- the GN slot}
\label{sec:approximations}

\textbf{Math.} $\Eqx[\bar f^\top \Sigma^{-1}\bar f]$. Apply Price's
theorem and the GN move (framework above): the $\partial_S$ derivative
is the residual-Hessian decomposition \eqref{eq:gn-hessian}, with the
sign-indefinite residual piece dropped and the covariance further
dropped, yielding \eqref{eq:Eff-lin}. This is the load-bearing
PSD-by-construction slot of the natural-grad. The exact
spectral-autocorrelation form \eqref{eq:gmix-exact-Eff} is the
alternative the GN framework rejected (sign-indefinite $\partial_S$).

\textbf{Value (code, line 297).} \code{trm += (del\_t)**2 * Eff},
where \code{Eff} $= \|\Eqx[\bar f]\|^2$ is stored as
\code{(Ef * Ef).sum(axis=-1)} (\code{drift\_moments\_gmix\_jax}
line 869). The forward holds only the $\Eqx[\bar f]^\top\Eqx[\bar f]$
piece of \eqref{eq:Eff-lin}; the $\tr$-piece is absorbed into the
gradient by the \code{custom\_vjp} below (saves a per-source
contraction). A closed form for the exact $\Eqx[\bar f^\top \bar f]$
exists --- \code{gmix\_E\_full\_Eff}
(\code{sing/efgp\_gmix\_qx\_moments.py:88}) --- but is \emph{not} called
by the production fit; only \code{precompute\_aux} from that module is,
and only to supply \code{GmixQxAux} to the V-Hutch path.

\textbf{Gradients.} \code{Eff} comes from
\code{eff\_with\_grads(m, S, eff\_const, ef\_const, jac\_const)}
(\code{efgp\_jax\_drift.py:57--75}); the \code{custom\_vjp} backward
returns
\begin{equation*}
\partial\code{Eff}/\partial m_t = 2g\,(\text{ef\_const}\cdot\text{jac\_const}),
\qquad
\partial\code{Eff}/\partial S_t = g\,(\text{jac\_const}^\top\text{jac\_const}),
\end{equation*}
with $g$ the upstream cotangent ($=(\Delta_i)^2$ from line 297),
\code{ef\_const}$=\Eqx[\bar f]$, \code{jac\_const}$=\Eqx[J_{\bar f}]$.
Reading these against the GN-truncated Price's identities:
\begin{itemize}\setlength\itemsep{2pt}
\item $\partial m_t$ injection $2g\,\Eqx[J_{\bar f}]^\top \Eqx[\bar f]$.
  By Price, $\partial_m \Eqx[\bar f^\top\Sigma^{-1}\bar f] =
  2\E[J^\top\Sigma^{-1}\bar f]$; production drops $\text{Cov}(J,\bar f)$
  giving $2\Eqx[J]^\top\Sigma^{-1}\Eqx[\bar f]$. \textbf{Covariance
  drop} (one of the GN approximations).
\item $\partial S_t$ injection $g\,\Eqx[J_{\bar f}]^\top \Eqx[J_{\bar f}]$.
  By Price, $\partial_S\Eqx[\bar f^\top\Sigma^{-1}\bar f]$ has the
  residual-Hessian decomposition \eqref{eq:gn-hessian}; GN drops the
  residual + $\text{Cov}(J,J)$, leaving
  $\Eqx[J]^\top\Sigma^{-1}\Eqx[J]\succeq 0$. \textbf{GN move on the
  load-bearing slot.}
\end{itemize}
Without the \code{custom\_vjp}, autodiff would return zero in both
slots (the forward \code{Eff} has no JAX-traced $(m, S)$-dependence).
Net contributions of (iii) to $\partial \code{trm}$:
\begin{equation*}
\partial m_t \;=\; 2(\Delta_i)^2\,\Eqx[J_{\bar f}]^\top\Eqx[\bar f],
\qquad
\partial S_t \;=\; (\Delta_i)^2\,\Eqx[J_{\bar f}]^\top\Eqx[J_{\bar f}]\;\succeq\;0.
\end{equation*}
\textbf{This is the GN slot:} $m$-gradient
from Part A$'$ (exact slot), $S$-gradient from Part B$'$ (Taylor's
trace piece). The Part B$'$ $\partial S_t$ is the load-bearing PSD slot
that protects SING's $\rho$-damped $J$-update.

\subsection*{Item (iv): the $V$ term --- production drops; App.~\ref{app:vhutch} opt-in via Price + Hessian-drop}

\textbf{Math.} $\Eqx[V(x_i)]$, where $V(x) =
\sum_r\sigma_r^{-2}\phi(x)^* A_r^{-1}\phi(x)$ is the $q(f)$ posterior
variance contribution. From
\begin{equation}
\Cov_{q(f)}\!\bigl[f_r(x), f_r(x')\bigr] = \phi(x)^* A_r^{-1}\phi(x'),
\label{eq:qf-cov}
\end{equation}
the full $q(f)$ second moment of $f$ splits as
\begin{equation}
\Eq[q(f)]\!\bigl[f(x)^\top \Sigma^{-1} f(x)\bigr]
\;=\; \sum_r \sigma_r^{-2}\bigl[\bar f_r(x)^2 + V_r(x)\bigr];
\label{eq:Eff-full}
\end{equation}
$V$ is the piece that (iii)'s $\bar f^\top \Sigma^{-1}\bar f$ doesn't
capture. $V$ has no natural Part A/B split --- it's a single integral.

\textbf{Production: value and gradients = 0.} The
$-\tfrac{\Delta_i}{2}\Eqx[V(x_i)]$ contribution is omitted from
\code{trm}. Geometric justification: translation-invariance of
stationary kernels gives $\nabla V(x) \approx 0$ wherever the source
set $\{m_i^{(k)}\}$ looks isotropic at scale $\ell$ around $x$; for
SDE-recovery with diverse ICs and ergodic mixing this holds at every
interior $m_i$, so the dropped $-\tfrac{\Delta_i}{2}\Eqx[\nabla V]$
contribution to the q(x) gradient is negligible. Fails at
boundary/edge points and cold-start.

\textbf{App.~\ref{app:vhutch} opt-in: Price + Hessian-drop.} Apply
Price's theorem \eqref{eq:price} to $g = V$:
\begin{equation}
\partial_{m_i}\Eqx[V] = \E[\nabla V],\qquad
\partial_{S_i}\Eqx[V] = \tfrac{1}{2}\E[\nabla^2 V],
\end{equation}
and apply the same Hessian-drop the framework uses elsewhere
($\E[\nabla^2 V] \to 0$). and drop \emph{only} the $\partial_S$ Hessian, keeping the value and
$m$-gradient exact:
\begin{equation}
\Eqx[V]\ \text{(exact)},\qquad
\partial_{m_i}\Eqx[V] = \Eqx[\nabla V]\ \text{(exact)},\qquad
\partial_{S_i}\Eqx[V] = \tfrac12\Eqx[\nabla^2 V]\to 0.
\label{eq:V-grads-taylor}
\end{equation}
This is exactly the item-(ii) $\bar f$ treatment: $\Eqx[V]$ and
$\partial_{m_i}\Eqx[V]$ are the \emph{same} Gaussian-averaged NUFFT (the
closed form \eqref{eq:V-hutch-closed}), computed off-graph and
differentiated through the $S$-envelope by autodiff --- so the injected
$\partial_m$ is the exact envelope-damped $\Eqx[\nabla V]$, \emph{not} the
point value $\nabla V(m_i)$ --- and only $\partial_S$ is Hessian-dropped.
$\omega_r$ ($=$ the $D_\theta A_r^{-1}D_\theta$ diagonal sums) is
Hutchinson-estimated; \code{FrozenEFGPDriftWithVHutch}'s
\code{custom\_vjp} injects the value and $m$-gradient and zeroes the
$S$-slot.

\subsection*{End-to-end: from \code{trm} gradients to a natural-grad step}

\begin{enumerate}\setlength\itemsep{2pt}
\item Pre-compute drift moments
  $(\code{Ef}, \code{Eff}, \code{Edfdx}) = (\Eqx[\bar f],
  \|\Eqx[\bar f]\|^2, \Eqx[J_{\bar f}])$ at all transition base points
  via \code{drift\_moments\_gmix\_jax}; stored on
  \code{FrozenEFGPDrift}.
\item \code{nat\_grad\_transition} (\code{sing/sing.py}) runs
  \code{vmap(grad(compute\_neg\_CE\_single))} over transitions.
\item \code{compute\_neg\_CE\_single} consumes \emph{mean parameters}
  $\mu = (\E[x_i], \E[x_i x_i^\top], \E[x_i x_{i+1}^\top])$, internally
  runs \code{mean\_to\_marginal\_params} to recover $(m_i, S_i, S_{i,i+1})$,
  and assembles \code{trm} from items (i)-(iii).
\item \code{jax.grad} chain-rules:
  (i) pure autodiff through traces;
  (ii) autodiff $+$ \code{ef\_with\_jac\_grad}'s injected
  $\partial$\code{Ef}/$\partial m$;
  (iii) entirely via \code{eff\_with\_grads}'s injected
  $\partial$\code{Eff}/$\partial(m, S)$.
  Returns $\partial\mathcal{L}/\partial\mu$ --- which by
  \eqref{eq:natgrad-duality} \emph{is} the natural gradient on
  $(J, h, L)$. SING $\rho$-damps and applies.
\end{enumerate}

% ============================================================================
\section{Algorithm summary and complexity}
\label{sec:algorithm}
% ============================================================================

\subsection*{One outer iteration of the E-step}

Given current $q(x_{0:T})$ (marginals $m_i, S_i$ and cross-covariances
$S_{i,i+1}$), kernel hyperparameters $\theta$, diffusion $\Sigma$, and
spectral grid $\{\xi_k\}$ with weights $D_\theta$:

\begin{enumerate}
\item \textbf{Increments.} $d_i = m_{i+1}-m_i$,\quad $C_i = S_{i,i+1}-S_i$.
\item \textbf{Sources.} \code{\_flatten\_stein} packs
  $\{(m_i,S_i,d_i,C_i,\Delta_i)\}_{k,i}$ into $N=K(T{-}1)$ flat
  supersources (padded slots get $S_i{=}I$ and weight $0$). For
  \code{'gmix'}/\code{'analytic'} the sources are the
  $q(x_i)=\mathcal{N}(m_i,S_i)$ Gaussians themselves; for \code{'mc'},
  draw $x_i^{(s)}\sim q(x_i)$.
\item \textbf{Precision.} Build the BTTB generator $T_r$ from weighted
  sources by the resolved backend (§\ref{sec:compute}: closed-form
  spreader / Taylor-envelope NUFFTs / plain Type-1 NUFFT) with weights
  $\Delta_i/\sigma^2$; $A_r$ acts as
  $v \mapsto v + D_\theta T_r D_\theta v$.
\item \textbf{RHS.} Same primitive with weights $d_{i,r}/\sigma^2$ and
  $[C_i]_{j,r}/\sigma^2$ produces $h_{1,r}$ and $h_{2,r}$.
\item \textbf{Solve.} Jacobi-preconditioned CG: $A_r\mu_r = h_r$.
\item \textbf{Drift moments.} All three backends first return point-eval
  moments (\code{drift\_moments\_jax}); with the default
  \code{qx\_moments\_method='gmix\_batched'} these are then
  \emph{overwritten} by the Gaussian-averaged
  $\Eqx[\bar f], \Eqx[J_{\bar f}], \|\Eqx[\bar f]\|^2$ from gmix-gather
  (\code{drift\_moments\_gmix\_jax}; §\ref{sec:moments}). Setting
  \code{'linearised\_shim'} keeps the point-evals (App.~\ref{app:impl}).
\item \textbf{$q(x)$ update.} Feed these into
  \code{compute\_neg\_CE\_single} via \code{FrozenEFGPDrift};
  \code{jax.grad} inside \code{nat\_grad\_transition} returns
  $\partial\mathcal{L}/\partial\mu$, which by
  \eqref{eq:natgrad-duality} is the natural gradient on $(J, h, L)$.
  SING $\rho$-damps and applies.
\end{enumerate}

\subsection*{Complexity}
\label{sec:complexity}

Let $N$ be the source count ($N = K(T{-}1)$ for the gmix path,
$N=K(T{-}1)S$ for MC; $K$ is the number of trials), $M$ the number of
Fourier features, $k$ the CG iteration count. The $q(f)$ update costs
\begin{equation}
\underbrace{O(N)}_{\text{spread / cloud assembly}}
\;+\;
\underbrace{O(M\log M)}_{\text{NUFFT / FFT}}
\;+\;
\underbrace{O(k\, D_x\, M\log M)}_{\text{CG over output dims}}.
\end{equation}
For \code{'analytic'} the $O(N)$ term becomes $O(P(p)N)$ with
$P(p)=\binom{p+2}{2}$, but as one batched transform. Drift moments at
$T$ points cost $O(T) + O(D_x M\log M)$. The optional $V$-restoration of
App.~\ref{app:vhutch} adds, \emph{per outer iter},
$O(P\,D_x\,k\,M\log M)$ for the $P$ Hutchinson probes that build
$\hat\omega_r$, and \emph{per inner iter} $O((1{+}D_x)(M\log M + N))$ for
the $1{+}D_x$ batched NUFFTs that produce $(V_i,\nabla V_i)$ at the
current $\{m_i\}$ (\code{efgp\_em.py} calls
\code{precompute\_V\_and\_grad\_E\_V\_per\_t} inside \code{one\_iter});
only the CG solves are hoisted to the outer iter. The custom-VJP shim
itself is a table lookup. The implementation rule is to never materialise
$K$, $\Phi$, $\Phi^*W\Phi$, or $A_r^{-1}$: every required action is a
Fourier matvec, a BTTB matvec, or a CG solve.

% ============================================================================
\appendix
% ============================================================================

% ============================================================================
\section{Math-to-code dictionary}
\label{app:dictionary}
% ============================================================================

\noindent\footnotesize
\begin{tabular}{@{}p{0.44\textwidth} p{0.52\textwidth}@{}}
\hline
Math & Code \\
\hline
Spectral grid weights $D_\theta = \diag(\sqrt{S(\xi_k)\,h^d})$
& \code{grid.ws} (real part) \\
Frequency points $\xi_k$
& \code{grid.xis\_flat} \\
Adjoint Type-1 NUFFT $F_X^{*}$
& \code{nufft1} \\
Forward Type-2 NUFFT $\phi(x)^*\mu$
& \code{nufft2} \\
BTTB generator $T_r = F_X^{*}W_r F_X$
& \code{bttb\_conv\_vec\_weighted} $\to$ \code{make\_toeplitz} \\
Operator $A_r v = v + D_\theta T_r D_\theta v$
& \code{make\_A\_apply} \\
Solve $A_r \mu_r = h_r$ (Jacobi-preconditioned CG)
& \code{cg\_solve} \\
$h_r = h_{1,r} + h_{2,r}$ build, \code{'mc'} backend
& \code{compute\_mu\_r\_jax} \\
same, \code{'gmix'} backend (§\ref{sec:gmix})
& \code{compute\_mu\_r\_gmix\_jax} $+$ \code{sing/efgp\_gmix\_spreader.py} \\
same, \code{'analytic'} backend (§\ref{sec:analytic})
& \code{compute\_mu\_r\_analytic\_jax}, \code{\_env\_taylor\_pairs} \\
$q(f)$ update $+$ drift moments in one call (per backend)
& \code{qf\_and\_moments\_\{jax, gmix\_jax, analytic\_jax\}} \\
$\bar f(m_i)$, $J_{\bar f}(m_i)$ point-eval (\code{'linearised\_shim'})
& \code{drift\_moments\_jax} \\
$\Eqx[\bar f]$, $\Eqx[J_{\bar f}]$ Gaussian-averaged (\textbf{default})
& \code{drift\_moments\_gmix\_jax} $+$ \code{sing/efgp\_gmix\_gather.py} \\
$\Eqx[\bar f^\top\bar f]$ exact closed form (\emph{defined, unused in the fit})
& \code{gmix\_E\_full\_Eff} (\code{sing/efgp\_gmix\_qx\_moments.py}) \\
Spectral autocorrelation $\rho(\delta)$ precompute (one FFT/outer iter)
& \code{precompute\_aux} (only feeds the V-Hutch path) \\
Multi-trial flatten of $\{(m_i,S_i,d_i,C_i,\Delta_i)\}_{k,i}$
& \code{\_flatten\_stein} \\
$\diag(\Phi A_r^{-1}\Phi^*)$ (general Hutchinson primitive; tests only)
& \code{hutchinson\_diag} \\
$\omega_r(\delta)$ Hutchinson estimator (eq.~\ref{eq:omega-hutch})
& \code{precompute\_omega\_per\_r} (\code{sing/efgp\_qx\_v\_hutch.py}) \\
$(V_i, \nabla V_i)$ batched NUFFT-2 precompute (eq.~\ref{eq:V-hutch-closed})
& \code{precompute\_V\_and\_grad\_E\_V\_per\_t} \\
V-Hutch custom-VJP shim into per-transition $q(x)$ update
& \code{FrozenEFGPDriftWithVHutch} (\code{efgp\_qx\_v\_hutch.py:485}) \\
GN custom-VJP shim into $q(x)$ update (\textbf{default}, §\ref{sec:interface})
& \code{FrozenEFGPDrift}, \code{ef\_with\_jac\_grad}, \code{eff\_with\_grads} \\
$q(x)$ consumer: \code{vmap(grad(trm))} $\to$ $(J,h,L)$
& \code{nat\_grad\_transition} (\code{sing/sing.py:46}) \\
\code{trm} assembly from (\code{Ef}, \code{Eff}, \code{Edfdx})
& \code{compute\_neg\_CE\_single} (\code{sing\_helpers.py:289}) \\
EM driver, backend resolution, M-step
& \code{fit\_efgp\_sing\_jax} (\code{efgp\_em.py:404}) \\
\hline
\end{tabular}

\medskip
\noindent
\begin{tabular}{ll}
\hline
$q(x)$ statistic & Code \\
\hline
$m_i$ (marginal mean) & \code{ms} \\
$S_i$ (marginal covariance) & \code{Ss} \\
$\Cov(x_i,x_{i+1}) = S_{i,i+1}$ & \code{SSs.transpose(0,2,1)} \\
$d_i = m_{i+1}-m_i$ & \code{d\_a} \\
$C_i = S_{i,i+1}-S_i$ & \code{C\_a} \\
\hline
\end{tabular}

% ============================================================================
\section{V-Hutch restoration (opt-in)}
\label{app:vhutch}
% ============================================================================

Default production drops the $V$ term entirely (§\ref{sec:approximations}). For problem classes where the geometric
argument is suspect --- boundary or cluster-supported coverage,
anisotropic coverage, sharp turns near $\partial\Omega$, cold-start outer
iters before $q(x)$ has localised --- this appendix gives a Hutchinson +
NUFFT restoration whose $(m_i, S_i)$ gradient structure is exactly the
linear case of §\ref{sec:approximations}. Wiring guard
(\code{efgp\_em.py:151--153}): the restoration is active only for
\code{qx\_moments\_method} $\in$ (\code{'linearised\_shim'},
\code{'gmix\_batched'}) \emph{and} \code{estep\_method} $\in$
(\code{'gmix'}, \code{'analytic'}); with \code{'mc'} it is silently
inactive.

\subsection{Closed-form $\Eqx[V]$ via BTTB-diagonal sums of $D A^{-1} D$}

$V_r$ is a \emph{quadratic form in the Fourier features}, so --- exactly
like the mean-square $\bar f_r^2=\phi^*\mu_r\mu_r^*\phi$ of item (iii) ---
it is a sum of difference-frequency modes,
\begin{equation}
V_r(x)=\sigma_r^{-2}\phi(x)^*A_r^{-1}\phi(x)
=\sigma_r^{-2}\sum_{k,\ell}D_k\,(A_r^{-1})_{k\ell}\,D_\ell\;
 e^{2\pi i(\xi_\ell-\xi_k)^\top x}.
\end{equation}
The dense $A_r^{-1}$ (unlike the rank-one $\mu_r\mu_r^*$) only sets the mode
coefficients; the $x$-dependence is still $e^{2\pi i\delta^\top x}$,
$\delta:=\xi_\ell-\xi_k$, so each mode averages under $q(x_i)$ by the same
characteristic function \eqref{eq:Eqx-chi}. Collecting modes by lag, define
the diagonal sums
\begin{equation}
\omega_r(\delta) \;:=\; \sum_{(k,\ell):\,\xi_\ell - \xi_k = \delta}
   D_k\,(A_r^{-1})_{k\ell}\,D_\ell,
\label{eq:omega-def}
\end{equation}
the BTTB projection of $D_\theta A_r^{-1} D_\theta$ onto lags $\delta$
on the $2M$-padded spectral-difference grid. By stationarity of the
kernel, $\omega_r(-\delta) = \omega_r(\delta)^*$ (Hermitian-symmetric;
generally complex when source locations break centrosymmetry). Then
\begin{equation}
\E_{q(x_i)}[V_r(x_i)]
\;=\; \sum_\delta \omega_r(\delta)\,
       e^{2\pi i\,\delta^\top m_i}\,
       e^{-2\pi^2 \delta^\top S_i \delta},
\label{eq:V-hutch-closed}
\end{equation}
a Gaussian-enveloped Type-2 NUFFT against the difference-grid coefficient
block --- the primitive \code{gmix\_E\_full\_Eff} already uses for
$\Eqx[\bar f^\top\bar f]$, with $\omega_r$ replacing the autocorrelation
$\rho_r$ of $\mu_r$.

So the recipe is modular: (i) obtain the drift posterior variance --- i.e.\
the coefficients $\omega_r(\delta)$ --- by \emph{any} estimator (the
Hutchinson estimator below, or a pathwise/Matheron posterior sampler,
Wilson et al.\ 2020); (ii) apply the gmix envelope-NUFFT
\eqref{eq:V-hutch-closed}. Only step (i) touches $A_r^{-1}$; step (ii),
the $q(x_i)$-average, is identical to every other moment.

\subsection{Hutchinson estimator of $\omega_r$}

$\omega_r$ is estimated unbiasedly by a padded-FFT pair on $P$ Rademacher
probes $z_p\!\in\!\{\pm 1\}^M$:
\begin{equation}
\hat\omega_r \;=\; \frac{1}{P}\sum_{p=1}^{P}\,
   \mathcal{F}^{-1}\!\bigl(\overline{\mathcal{F}[D u_{p,r}]}
   \cdot \mathcal{F}[D z_p]\bigr),\quad
u_{p,r} = A_r^{-1} z_p,
\label{eq:omega-hutch}
\end{equation}
with $A_r^{-1} z_p$ via the same Jacobi-preconditioned CG that drives
the $q(f)$ update. Variance scales as $\|D A^{-1} D\|_F/\sqrt{P}$;
$P\!=\!4$ probes shared across the $D_{\rm out}$ output dims are
sufficient. \emph{Keep the imaginary part:} a \code{.real} cast on the
IFFT discards $-\Im\omega(\delta)\sin(2\pi\delta^\top m)$, the dominant
error vs dense ground truth when sources are not centrosymmetric.

\subsection{Per-transition gradient: exact in $m$, $S$-slot dropped}

Following \eqref{eq:V-grads-taylor}, the shim keeps the value and
$m$-gradient exact and drops only the $S$-slot:
\begin{equation}
\E_{q(x_i)}[V(x_i)]\ \text{(exact)},
\qquad
\partial_{m_i}\E[V] = \E[\nabla V]\ \text{(exact)},
\qquad
\partial_{S_i}\E[V] = 0.
\end{equation}
\code{precompute\_V\_and\_grad\_E\_V\_per\_t} evaluates the closed form
\eqref{eq:V-hutch-closed} at $\{m_i\}$ and takes \code{jax.grad} w.r.t.\
$\{m_i\}$, so $\partial_m$ is the exact envelope-damped $\E[\nabla V]$
(not $\nabla V(m_i)$); \code{FrozenEFGPDriftWithVHutch} zeroes the
$\partial V_i/\partial S_i$ slot. This is the V-analog of the item-(ii)
$\bar f$ shim, at the same truncation (drop $\partial_S$ only).

The value $V_i$ and the $m$-gradient are exact (the closed form
\eqref{eq:V-hutch-closed} and its autodiff); the \emph{only} dropped term
is the $S$-slot $\partial_{S_i}\E[V]=\tfrac12\Eqx[\nabla^2 V]$, of size
$O(\|S_i\|\cdot\sigma_f^2/\ell^2)$ for SE kernels --- small for $q(x)$
localised at scale $\sim\ell$. To restore this $S$-pull exactly, replace
$V$ by its second-order Taylor with PSD-projected Hessian
$H_V^{+}\succeq 0$ (Price's identity, $\partial_S = \tfrac12 H_V^{+}$);
cost: $D(D{+}1)/2$ extra batched NUFFTs against frequency-modulated
$\omega(\delta)\cdot(-2\pi^2\delta_j\delta_k)\,\mathrm{env}(\delta)$
coefficient grids. Not enabled by default.

\subsection{Homogeneous-$S$ vs heterogeneous-$S$ sub-modes}

Two production sub-modes for the per-source envelope in
\eqref{eq:V-hutch-closed}:

\paragraph{\code{'hutch'} (default \emph{when} restoration is on;
\code{'none'} is the production default).}
$S_i$ is replaced by the trajectory-mean $\bar S$ in the envelope
$e^{-2\pi^2 \delta^\top S \delta}$, so the envelope folds into the
coefficient grid and the per-source NUFFT-2 is a standard (non-gmix)
transform. $\hat\omega_r$ is built once per \emph{outer} iter;
$(V_i, \nabla V_i)$ are rebuilt each \emph{inner} iter at the current
$\{m_i\}$ (and $\bar S$ from the current $S$'s).

\paragraph{\code{'hutch\_hetS'}.}
Per-source $S_i$ envelope, implemented via the Type-2 gmix-gather
primitive (\code{efgp\_gmix\_gather.py}). Pipeline: (i) inverse-FFT
$\hat\omega(\delta)$ once to a uniform spatial grid
$\widetilde\omega(x)$; (ii) for each source $i$, gather $\widetilde\omega$
via a per-source Gaussian stencil weighted by $\mathcal{N}(x; m_i, S_i)$
and Riemann-summed over a $(2r{+}1)^D$ neighbourhood, with
$r = \lceil n_\sigma\,\sigma_{\max}/h_{\rm grid}\rceil$ (static, picked at
startup from $V_0$). Both $V_i$ and $\nabla V_i$ flow through the same
primitive (gradient via \code{jax.grad} reverse-mode through the gather),
rebuilt every inner iter.

% ============================================================================
\section{Implementation and benchmarks}
\label{app:impl}
% ============================================================================

\subsection{Custom-VJP shims: why a shim, not pure autodiff}

\code{nat\_grad\_transition} runs
\code{vmap(grad(loss\_one\_transition))} where
\code{loss\_one\_transition} closes over a single $(m_t, S_t)$. Inside
that vmap, JAX/XLA does \emph{not} hoist the FFT across vmap'd
\code{nufft2} calls: a per-source NUFFT-2 inside \code{ff(t, m, S)}
pays the full $O(M\log M)$ FFT cost \emph{per source}, giving
$O(NM\log M)$ per inner iter and ${\sim}50\times$ slowdown at
$K{=}10/T{=}500$ (measured). \code{FrozenEFGPDrift} sidesteps this by
precomputing $(\bar f, J_{\bar f})$ at all sources via ONE forward
batched NUFFT-2 plus its AD-driven adjoint ($D$ Type-1 NUFFTs)
\emph{outside} the per-transition vmap, then injecting both as constants
per-transition via \code{custom\_vjp} with the §\ref{sec:qx} gradient
table. The V-Hutch shim
\code{FrozenEFGPDriftWithVHutch} follows the same pattern for
$(V_i, \nabla V_i)$. A refactor of \code{nat\_grad\_transition} to
compute the multi-transition ELBO in one shot and \code{jax.grad} it
once over all $(m_i, S_i)$ would recover the same asymptotic cost
without shims.

\subsection{Point-eval alternative: linearised-shim}

The §\ref{sec:moments} default Gaussian-averages $(\bar f, J_{\bar f})$
under $q(x_i)$ via the gmix-gather primitive. A cheaper alternative,
\emph{linearised-shim}, uses point-eval $(\bar f(m_i), J_{\bar f}(m_i))$
directly (no Gaussian averaging): one batched Type-2 NUFFT on $\mu_r$
at the $\{m_i\}$ source set, with no gather. The Taylor-truncated
gradient structure of §\ref{sec:qx} is identical --- so both consume
the same \code{FrozenEFGPDrift} custom-VJP --- but the base values
miss the first-order $S_i$-correction. Empirically (§\ref{app:impl}
wall-time table) the two paths reach the same fixed point to 3--4
decimals and within ${\sim}7\%$ wall, so gmix-batched is the production
default but linearised-shim is retained for problems where the
gather truncation is not negligible (small $\ell$ regime) or for
debugging.

\subsection{Wall-time comparisons ($K=10$, $T=500$)}

Benchmarks against the \code{'gmix\_batched'} baseline,
subprocess-isolated timings, warm JIT:

\begin{itemize}
\setlength\itemsep{2pt}
\item Base-point choice: \code{'gmix\_batched'} reaches the same
  $(\hat\ell, \hat\sigma_f^2, \text{lat\_rmse})$ fixed point as
  \code{'linearised\_shim'} to 3--4 decimals, at wall within ${\sim}7\%$
  (\code{demos/bench\_qx\_moments\_modes.py}).
\item V-restoration overhead: \code{'hutch'} ${\sim}53.7$\,s warm vs
  \code{'none'} ${\sim}48.3$\,s (${\sim}11\%$);
  \code{'hutch\_hetS'} ${\sim}64.8$\,s (${\sim}1.34{\times}$).
  Recovery is bit-identical across all three modes
  (\code{demos/bench\_hutch\_vs\_none.py}). \emph{Caveat:} the
  \code{efgp\_em.py} docstring quotes ${\sim}22\%$ for
  \code{'hutch\_hetS'} at the same $K,T$ --- the two figures come from
  different runs; re-measure before quoting either in the manuscript.
\end{itemize}

\subsection{$\nabla V$ magnitude diagnostic}

Setup: $K{=}10$ trials, $x_0\sim\mathcal{U}([-2,2]^2)$, integrated for
$8$\,s with $\sigma{=}0.4$, $\ell_{\rm true}{=}0.6$. Latent paths cover
$\Omega\approx[-3.5, 3]^2$ with $\sim$5{,}000 sources,
$|\Omega|/\ell^d \approx 120$; the §\ref{sec:approximations}-A
isotropic-coverage argument applies with margin.

Direct measurement at a representative transition shows the wiring is
correct (the \code{custom\_vjp} bwd produces $\nabla V$ matching the
precomputed \code{grad\_V\_const} to numerical precision) but the
\emph{magnitude} of $\nabla V$ is ${\sim}5\times 10^{-5}$ of the
$\partial \mathrm{Eff}/\partial m = 2 J^\top \E[\bar f]$ term per
transition, translating to ${\sim}0.16\%$ change in the full
\code{compute\_neg\_CE\_single} gradient direction. SING's $\rho$-damping
washes this out. The bit-identical recovery between \code{'none'} and
\code{'hutch'} (final $\hat\ell{=}1.1337$, $\hat\sigma_f^2{=}1.3230$,
latent RMSE${=}0.0646$, all matching to 3--4 decimals across all three
V-restoration modes) reflects the structural argument, not a wiring bug.

\subsection{Op count for V-Hutch per outer iter}

{\footnotesize
\begin{tabular}{lcc}
\hline
Operation & Count & Per-call cost \\
\hline
Hutch CG ($A_r^{-1} z_p$, eq.~\ref{eq:omega-hutch})
   & $P\cdot D_{\rm out}$
   & $O(k_{\rm cg}\, M\log M)$ \\
FFT cross-correlation (gives $\hat\omega_r$)
   & $P\cdot D_{\rm out}$
   & $O(M\log M)$ \\
Forward batched NUFFT-2 (gives $\{V_i\}$) \emph{---per inner iter}
   & 1
   & $O(K(M\log M + T))$ \\
Adjoint NUFFT-1 ($\nabla V_i$ via \code{jax.grad}) \emph{---per inner iter}
   & $D$
   & $O(K(M\log M + T))$ \\
\hline
\multicolumn{2}{l}{Per-inner-iter CG solves added}
   & 0 \\
\hline
\end{tabular}}

\medskip
NUFFTs are the only ops that scale with $N=K(T{-}1)$. Only the
$\hat\omega_r$ block (CG $+$ FFT cross-correlation) is hoisted to the
outer iter; the $1{+}D$ batched NUFFTs that produce $(V_i,\nabla V_i)$
sit \emph{inside} \code{one\_iter}, since they must be re-evaluated at
the current $\{m_i\}$, so they scale with $n_{\rm estep}$. At
$K{=}10/T{=}500$ ($D{=}2$, $M\!\sim\!10^3$, $P{=}4$,
$k_{\rm cg}\!\sim\!30$): CG block ${\sim}0.1$\,s per outer iter, NUFFT
block ${\sim}1$\,ms per inner iter --- consistent with the
${\sim}11\%$ measured wall for \code{'hutch'}.

\end{document}
